\chapter{Criminal Proceedings}


% \section{Nature of Criminal Proceeding}
% \begin{enumerate}[label=$(\text{\roman*})$, itemsep=1em]
%    \item \textbf{Adversarial Process:}
%    In Bangladesh, the criminal justice system follows an adversarial process, where the trial is a contest between the state (prosecution) and the accused. The court remains neutral, does not prepare the case, and only decide within a complex set of rules, whether the accused is proven guilty.
   
%    \item \textbf{Presumption of Innocence/ Criminal Standard of Proof:}
%    A person accused of a crime is presumed to be innocent until the prosecution proves his guilt beyond every reasonable doubt if there is a little doubt in proving the elements of th offense concerned, the accused will be set free.

%    \item Criminal justice punishes wrongdoing, not enforces rights. The victim accuses the defendant but doesn't claim a right. The court does not enforce duties or rights, but penalizes for duties already broken or rights already violated. In serious offenses, the state prosecutes the offender, even without a separate accuser.

%    \item All criminal proceedings in Bangladesh are regulated under the Code of Criminal Procedure Code 1898 and Criminal Rules and Orders unless otherwise excluded or specifically provided for.

%    \item Four agencies are involved in a criminal administration of justice: Police, Prosecution, Courts, Jail and Probation authority.
% \end{enumerate}


\section{Stages of Criminal Proceedings}

\subsection{Pre-Proceeding Stage}
\begin{enumerate}[label=$({\roman*})$, itemsep=1em]
   \item \textbf{FIR in Cognizable Offense:}
   \begin{itemize}
      \item Any person can file an FIR under section 154 for a cognizable offense.
      \item Police may also get info under section 157 from other sources.\item Magistrate may send a case to police for investigation under sections 190, 155(3), 156(3)
      % \item or refer it for inquiry under section 202 after receiving a complaint under section 200. The officer-in-charge must record the info, which becomes the FIR.
      \item Magistrate may accept a complaint under section 200 and send for inquiry under section 202.
   \end{itemize}
      
   \item \textbf{Complaint in Non-Cognizable and Cognizable Offense:}
   % For a non-cognizable offence, police cannot investigate without a Magistrate's order; they record the info in the General Diary and refer the informant to the Magistrate, who examines him under oath (sections 155 \& 200). Any person may also complain in writing to a Magistrate, who upon examination may take cognizance (C.R. Case), dismiss it, or order inquiry/investigation (C.R.P. Case).
   Police cannot investigate non-cognizable offences without Magistrate's order; they record in General Diary and forward informant. Magistrate examines under oath (sections 155, 200) and may take cognizance (C.R. Case), dismiss, or order inquiry/investigation before cognizance (C.R.P. Case).

   \item \textbf{Reporting to the Magistrate:}
   FIR filed at the police station gets FIR No./P.S. Then the original copy must be sent to the Magistrate within 24 hours (section 157, CrPC). Once received, the case gets a G.R. Case No. which is purely administrative and automatic.
   % In complaint cases, no P.S. Case exists; the matter starts as a C.R.P. Case, and if the Magistrate sees reason, it becomes a C.R. Case.
   % FIR filed in police station gets FIR No./P.S. Case No. Almost all criminal cases except complaints start this way. Original FIR sent to Magistrate within 24 hrs (section 157, CrPC). On receipt, Magistrate court assigns a G.R. Case number, which is purely administrative and automatic.
   Complaint cases start with C.R.P. or C.R. Case number in Magistrate court without a P.S. Case number. If Magistrate finds reason to proceed, the petition case becomes a C.R. Case.
   
   \item \textbf{Investigation and maintaining Case Diary:}
   For a cognizable offense, police may investigate without Magistrate's order (section 156). But for a non-cognizable offense, police can investigate only after Magistrate's order (section 155).

   \item \textbf{Final Report/ Charge Sheet:}
   After investigation, police may submit either a Final Report or Charge Sheet without court interference (section 173). A Final Report means no offense is made out and the accused may be released. A Charge Sheet recommends prosecution. The report is sent to the Magistrate empowered to take cognizance.

   \item \textbf{Final Report and Naraji Petition}
   If police submits a final report with no case, the Magistrate may accept or reject it. On rejection, he may order further investigation or take cognizance after examining the complainant. If accepted, the informant may file a Naraji Petition, treated as a complaint, and the Magistrate may issue process or order inquiry by another Magistrate.
\end{enumerate}

\subsection{Proceeding Stage}
\begin{enumerate}[label=$({\roman*})$, itemsep=1em]
   \item \textbf{Taking Cognizance:} it is the condition precedent to the initiation of proceedings. Under section 190, any CMM, CJM, First Class Magistrate, or specially empowered Magistrate may take cognizance from $(i)$ police report, $(ii)$ complaint, or $(iii)$ own knowledge/private information.
   \item \textbf{How Cognizance is taken:} Normally in front of the accused. If in custody, accused is brought to court; if not arrested, Magistrate issues summon or warrant. Magistrate examines case record to see if judicial proceeding can be initiated. For complaint cases, cognizance is taken in absence of accused as per section 200.
   \item \textbf{Start of a Criminal Proceeding:} If cognizance is taken on a charge sheet, criminal proceeding starts at once. Otherwise if taken on the basis of complaint/private info, it starts only after inquiry/investigation and issue of process under section 204. For non-cognizable offences, police record in General Diary, forward the informant; Magistrate examines the informant under oath (section 200), then may order inquiry/investigation. Proceeding begins once process (summons/warrant) is issued.
\end{enumerate}



\subsection{Trial Stage}
\subsubsection{Trial in the Magistrate Court}

\begin{enumerate}[label=$({\roman*})$, itemsep=1em]
   \item \textbf{Pre-trial Hearing / Discharge Before Framing of Charge}
   Before trial begins, the accused must appear or be brought to court. If absent while on bail, a bench warrant may be issued. The Magistrate reviews the case, and if the charge seems groundless, the accused may be discharged under section 241A. This is called "no case to answer" when the prosecution fails to show a prima facie case.
   
   \item \textbf{Framing of Charges:}
   If the Magistrate finds a prima facie case and has jurisdiction, he shall frame a formal charge (section 242). The charge must state time, place, person, and details so the accused knows the allegation. By this formal trial begins.
   
   \item \textbf{Plea and Conviction:}
   After the charge is framed, the accused is asked if he pleads guilty. If he admits the offence, the Magistrate may convict him under section 243.
   
   \item \textbf{Hearing/ Taking Evidence:}
   If the accused doesn't plead guilty, the Magistrate will hear the case based on evidence. The accused and witnesses are examined and cross-examined as per law (section 244).
   
   \item \textbf{Examination of the Accused:}
   After prosecution witnesses are examined, the court may question the accused under section 342 to explain points from the evidence. The accused may answer or remain silent, and the court may draw inferences from the response or refusal.
   
   \item \textbf{Acquittall:}
   If after hearing evidence and examining the accused the Magistrate finds the accused not guilty, he shall record an order of acquittal.
   
   \item \textbf{Sentence and transfer for Sentence:}
   If the Magistrate finds the accused guilty, he will pass sentence. But if the punishment needed is more severe than he can give, he will forward the case to the CJM or First Class Magistrate under section 349.
\end{enumerate}

\subsubsection{Trial Stage in Sessions Court}

\begin{enumerate}[label=$({\roman*})$, itemsep=1em]
   \item \textbf{Opening of the Prosecution Case:}
   When the accused appears, the Public Prosecutor opens the case by describing the charge and outlining the evidence to prove guilt (section 265B).
   
   \item \textbf{Pre-trial Hearing/ Discharge before framing of Charge:}
   After the prosecution opens the case, the Session Judge allows both sides to argue for framing charge or discharge (no witness examination yet). If the Judge finds no prima facie case, he discharges the accused with reasons (section 265C). This is when the defense can raise “no case to answer” before charges are framed.
   
   \item \textbf{Framing of Charge:}
   If, on the other hand, it shall frame a charge (section 265D). Formal trial starts with the framing of charge.
   
   \item \textbf{Plea and Conviction:}
   After framing the charge, it is read and explained to the accused, who is asked to plead guilty or claim trial. If guilty, the Court records the plea and may convict (section 265E).
   
   \item \textbf{Prosecution Evidence:}
   If the accused doesn't plead guilty or the plea is rejected, the court sets a date for examining witnesses. The Public Prosecutor first examines prosecution witnesses, followed by cross-examination by the accused and any re-examination (sections 265F and 265G).
   
   \item \textbf{Acquittal:}
   If the Court finds no evidence that the accused committed the offense after considering prosecution evidence and defense, it shall acquit the accused (section 265H).
   
   \pagebreak
   \item \textbf{Defense Evidence:}
   If the accused is not acquitted, he can present his defense and call witnesses. Defense witnesses are examined by the defense counsel, then cross-examined by the Public Prosecutor, followed by any re-examination (section 265I).
   
   \item \textbf{Summing Up/ Closing of the Prosecution and Defense Case:}
   After hearing evidence, the prosecutor sums up the strength of their case and weaknesses of the defense, explaining the burden of proof. Then the defense replies, highlighting their case's strength and prosecution's weaknesses to show the accused is not proved guilty beyond reasonable doubt (section 265J).
\end{enumerate}


\subsection{Post Trial Stage}
\begin{enumerate}[label=$(\roman*)$, itemsep=1em]
   \item If it is an acquittal and the accused is in jail, a copy of the judgment is sent to the jail authority, which releases the offender as per judgment.
   \item if it is a conviction of imprisonment, the accused is taken from court to jail to serve the sentence unless suspended or remitted by the Government.
   \item When a sentence is fully executed, the executing officer returns the warrant to the issuing court with endorsement certifying execution (section 400).
   \item For death sentence, proceedings go to the High Court Division; execution occurs only if confirmed (section 374).
   \item Persons under 16 may be confined in reformatory, training centre, or serve probation instead of criminal jail (section 399).
\end{enumerate}
